\section{Discussion and Limitations}
\label{sec:discussion}
\paragraph{Answer to the question.} The verdict is \emph{conditional, and mostly favourable} for the sampling noise itself. In the regime where Kernel SHAP and permutation SHAP are genuinely stochastic ($d=16$ and $d=20$) and at budgets from a quarter of the library default upward, one run identifies the top-3 features of a prediction, its signs and its overall magnitude profile reliably by the pre-stated rule. The reliability does not extend to (i) the library's default $\ell_1$ setting for $d>10$, which introduces a bias that repetition cannot reveal (this is a pitfall of one implementation, shown at one budget), (ii) the ordering of features whose attributions are close (on German Credit, only $\nGerLowPairsResMin$--$\nGerLowPairsResMax$ of the pairs at ranks four and below are resolved at twice the noise), and (iii) a minority of instances with markedly lower agreement between runs. For low-dimensional data the question is moot because the default is exact. A practitioner who cannot repeat the explainer should (a) switch off $\ell_1$ selection or check it against an exact or high-budget reference on a sample of instances, (b) read the few leading features rather than a full ranking, and (c) be aware that at a 50-row background the choice of background changed the attribution vector much more than the explainer seed did (the ratio depends on the background size and the budget, neither of which we varied). If model evaluation is fast, $16\times$ the default costs about $\nGerSixteenSec$\,s per explanation on German Credit on our shared host and lowers the already small error further (Figure~\ref{fig:budget}).

\paragraph{Limitations.} The claims are limited to the following. \emph{Few stochastic datasets:} two ($d=20$ with 1{,}000 rows; $d=16$ with a 40{,}000-row subsample of one state-year), so dataset-level uncertainty is not estimated; larger $d$ (e.g.\ hundreds of features) is not covered and plausibly changes the answer. \emph{Fixed design choices:} one background size (50 rows), one model fit per model class (refit variation is a 10-seed probe on 20 instances), small non-tuned model configurations, an interventional value function, one library (\texttt{shap} 0.51.0) and its implementation of Kernel SHAP and permutation SHAP; other estimators \citep{covert2021improving,musco2025leverage} are not evaluated. \emph{Targets:} accuracy is against exact targets for the same background; it is not accuracy with respect to any ground truth of ``importance'', and a background redraw changes the target. The German MLP accuracy rows use 20 instances, and the $16\times$ cell uses 20 instances on German Credit only. \emph{Instances and thresholds:} 50 instances per dataset, half chosen near the decision boundary; the decision-rule thresholds are conventions fixed before the runs (they are not tuned), and the fraction-over-threshold results should be read over their sweep. \emph{Compute:} the budgets are CPU-scale; wall-clock numbers come from a shared, oversubscribed host and are not benchmarks. \emph{Not run:} the background-size axis, comparisons with other explainers and higher-dimensional datasets (Give Me Some Credit, Arrhythmia) were cut for budget reasons and are not reported as done.

\paragraph{Broader impact.} Attributions support high-stakes decisions. Our results suggest that a single run is rarely the weakest link at realistic budgets, which should not be read as a reason to trust attributions in general: the conclusions concern estimator noise for a fixed model, background and value function, not whether the explanation is faithful, causal or fair.
